%\documentclass{article}
%\usepackage{graphicx}
%\usepackage{amsmath}
%\usepackage{icml2020}
%\twocolumn

%\begin{document}

\section{BERT Finetuning Hyperparameters}
\label{app:berthparams}

Table \ref{tab:bert_tasks_hparams} presents the hyperparameters used for each model and task during finetuning.

\begin{table}[hbt!]
\footnotesize
\begin{center}
\caption{Hyperparameters for finetuning BERT model on downstream tasks.}
\label{tab:bert_tasks_hparams}
\begin{tabular}{c|c|c|c|c} \hline \hline
Task      & Model & Batch & Learning & Training \\
          &       & size  & rate     & epochs   \\ \hline
          & 336M  &       &          &          \\ 
MNLI      & 1.3B  & 128   & 1e-5     & 10       \\ 
          & 3.8B  &       &          &          \\ \hline 
          & 336M  &   128    & 5e-5     &          \\ 
QQP       & 1.3B  & 128   & 3e-5     & 12      \\ 
          & 3.8B  &  256     & 4e-5     &          \\ \hline 
          & 336M  & 64    & 3e-5     &          \\ 
SQUAD 1.1 & 1.3B  & 48    & 3e-5     & 2        \\ 
          & 3.8B  & 48    & 1e-5     &          \\ \hline 
          & 336M  & 48    & 3e-5     &          \\ 
SQUAD 2.0 & 1.3B  & 64    & 3e-5     & 2        \\ 
          & 3.8B  & 48    & 1e-5     &          \\ \hline 
          & 336M  & 32    & 2e-5     &          \\ 
RACE      & 1.3B  & 16   & 1e-5     & 3        \\ 
          & 3.8B  & 32    & 2e-5     &          \\ \hline 
\end{tabular}
\end{center}
\end{table}


\section{Model Parallel Supplementary Material }
\label{sec:modelpar:supp}

In this section, we present further details about the hybrid model and data parallelism and handling random number generation.


\subsection{Hybrid Model and Data Parallelism}
Model parallelism is orthogonal to data parallelism, and so we can use both simultaneously to train large models in a reasonable amount of time. Figure \ref{fig:hybrid-mpdp} shows a grouping of GPUs for hybrid model and data parallelism. Two or more GPUs within the same server form model parallel groups (for example GPUs 1 to 8 in Figure \ref{fig:hybrid-mpdp}), and contain one instance of the model distributed across these GPUs.  The remaining GPUs, which could be within the same server but more typically are located in other servers, run additional model parallel groups. GPUs with the same position in each of the model parallel groups (for example GPUs 1, 9, ..., 505 in Figure \ref{fig:hybrid-mpdp}) form data parallel groups so that all GPUs within a data parallel group hold the same model parameters. During back propagation we run multiple gradient all-reduce operations in parallel to reduce weight gradients within each distinct data parallel group. The total number of required GPUs is the product of the number of model and data parallel groups. For example, for the 8.3 billion parameter model we use 8 GPUs per model parallel group and 64-way data parallelism, for a total of 512 GPUs.  All communication is implemented in PyTorch by Python calls to NCCL.  GPUs within each model parallel group perform all-reduces amongst all GPUs within the group.  For data parallelism, each of the all-reduce operations takes place with one of the GPUs from each model parallel group.


\subsection{Model Parallel Random Number Generation}
Techniques that utilize random number generation, such as dropout, are a staple of modern deep learning training. Transformers have dropout layers outside the model parallel regions before residual connections and within model parallel regions in the self attention block. Because some dropout layers are in a model parallel region, while others are not, we need to treat random number generation carefully to ensure dropout works correctly. To synchronize residual connection dropout across model parallel workers we seed the random number generators at the beginning of training with the same seed. This results in identical dropout patterns across all model parallel workers. However, dropout within a model parallel region should result in different random patterns for each worker to achieve randomness across the entire operation. To achieve this we maintain a separate random number generator for dropout within model parallel regions. This random number generator is uniquely seeded for each model parallel worker.

\begin{figure}[h]
\begin{center}
 \includegraphics[scale=0.32]{./hybrid_mpdp_5.png}
 \caption{Grouping of GPUs for hybrid model and data parallelism with 8-way model parallel and 64-way data parallel.}
 \label{fig:hybrid-mpdp}
\end{center}
\vspace{-6mm}
\end{figure}

\section{Text Samples}
\label{app:samples}
Below are some text samples generated by Megatron-LM using a context prompt. Some of the texts are cut short.
\fbox{
\parbox{0.48\textwidth}{
\rule{0.48\textwidth}{0.4pt}
{\bf Context}: There is a major force uniting America’s fiercely partisan politicians: big technology companies. Democrats and Republicans at the federal and state levels are coming together to scrutinize the power of the Silicon Valley giants and, potentially, to rein them in. Letitia James, the Democratic attorney general of New York, announced on Friday that attorneys general in eight states — four Democrats and four Republicans — and the District of Columbia had begun an antitrust investigation of Facebook.

\rule{0.48\textwidth}{0.4pt}
{\bf Megatron-LM}: The probe, which Ms. James said was likely to result in a lawsuit, is the latest in a string of high-profile regulatory actions, in state and federal courts, that challenge the market dominance of tech companies. The Justice Department is also investigating Facebook's data practices and whether the company violated a 2011 consent decree with the Federal Trade Commission that required the company to obtain explicit consent from users before sharing their personal information with third parties.
\\[1\baselineskip]
What might prompt a bipartisan coalition to sue Facebook? It isn't clear that there's any one thing. But lawmakers and regulators have been increasingly alarmed by the company's recent controversies, which include a data breach affecting more than 50 million users, the spread of fake news on the platform, and the spreading of propaganda by Russian-backed accounts during the 2016 election.
\\[1\baselineskip]
It's also not clear how serious the Democratic attorneys general are about filing a lawsuit. While the attorneys general announced their investigation in a press release, they said they had not yet reached a decision on whether to sue. In the past, some of the attorneys general have focused their efforts on trying to extract settlements from the tech companies without filing lawsuits.
\\[1\baselineskip]
"As attorneys general who have investigated and prosecuted health-care fraud, price-fixing, and other violations of our states' antitrust laws, we are committed to holding companies accountable when they fail to protect consumers and competition in the market," Ms. James said in a statement.
\\[1\baselineskip]
The investigation, which was first reported by The New York Times, was announced as the companies that dominate the tech industry are facing scrutiny from antitrust regulators around the world. Google, Twitter, Facebook and Amazon have faced a series of antitrust probes and class-action lawsuits in the United States and Europe ... 
%Last month, European regulators fined Google \$2.7 billion for unfairly promoting its own shopping service. And the European Union's competition commissioner, Margrethe Vestager, has threatened to hit Google with additional penalties.

}}

\fbox{%
\parbox{0.48\textwidth}{%
\rule{0.48\textwidth}{0.4pt}
{\bf Context}: A flower, sometimes known as a bloom or blossom, is the reproductive structure found in flowering plants (plants of the division Magnoliophyta, also called angiosperms). The biological function of a flower is to effect reproduction, usually by providing a mechanism for the union of sperm with eggs.

\rule{0.48\textwidth}{0.4pt}
{\bf Megatron-LM}: The structural components of the flower are the petals and sepals, which are the reproductive organs, the hypanthium, which forms the base of the flower, and the calyx, which is the stigma. A flower is a composite structure, made up of sepals and petals, the sepals being composed of three fused tepals and the petals of two fused petals, though there are a few exceptions to this rule. The three sepals are fused to form a tube called the calyx, which is the structure that attracts the pollinating insects. The petals are composed of two fused sepals (the outer ones) which open and fold inward (clasp) at the tips, forming the petals. The four sepals and petals are united to form a corolla, the outermost structure of a flower that carries pollen.
\\[1\baselineskip]
The structure of the flower has important functional significance, including protection of the ovary from insects and the transfer of pollen between the male and female parts of a flower. The structure of the flower is also used by the plants to attract pollinators, for example by the shape of the petals and the colours of the stamens.
\\[1\baselineskip]
The structure of a flower varies widely between plants. The simplest flowers are unbranched and made up of only two sepals and a single petal. Other examples are composed of several petals joined together, with the innermost sepals completely fused to form a tube called the calyx. The petals are usually coloured to attract insects, which can only see certain wavelengths of light. The sepals are often coloured too, but usually have a more limited range of colours.
\\[1\baselineskip]
Some flowers have additional structures, such as stamens, anthers and ovaries. The presence of stamens and anthers in some flowers is used for cross-pollination. The ovary is the organ where the female gametes are produced, and is the first structure that is developed in an embryo of a flower. The ovary is the female flower's reproductive system, and consists of the ovary wall, the ovule, and the oviduct.
\\[1\baselineskip]
Flowers vary in their size and shape, from tiny to huge, and are often brightly coloured. Flowers may be seen as part of a larger structure, such as a flower head, an inflorescence, a bract, a spathe, or a carpels.
\\[1\baselineskip]
The term "flower" is usually applied to the floral structures that are visible externally, but many ...
%flowers produce their flowers internally, and these structures are collectively known as the "floral anatomy" or "floral anatomy of a plant". The floral anatomy includes structures that are only visible internally, such as the ovary, the style, and the stigma.
%\newline \newline
%The term flower is also applied to the fruit of a flowering plant, such as the berry of a strawberry. The term may also be used to refer to the various parts
%of the flower that are not visible externally.
%The term "flower" is also used in other contexts, including the common names of certain plants and the eponymous genus of an entire family of flowering plants, Asteraceae.

}}

\fbox{%
\parbox{0.48\textwidth}{%
\rule{0.48\textwidth}{0.4pt}
{\bf Context}: OpenAI did not release the full GPT-2 model due to concerns of malicious use, but they did release a smaller version equivalent in size to the original GPT (117 M parameters), trained on the new, larger dataset. Although not as powerful as the large model, the smaller version still has some language generation chops.

\rule{0.48\textwidth}{0.4pt}
{\bf Megatron-LM}: In December 2016, DeepMind released an even larger model, GPT-3, with more than 1.5 billion parameters. DeepMind has no plans to release the model, so the question remains whether other AI researchers will be able to copy and build on it.
\\[1\baselineskip]
Other language models have been built using other techniques. In 2016, a team of researchers at Columbia University, the University of Toronto, and Google DeepMind showed that neural network-based models, including deep neural networks, can learn the grammatical rules of a language. A group at Google DeepMind demonstrated that these rules could be represented in a network as a kind of decision tree, where the nodes in the tree are the words in the sentence, and the rules for the words are represented by the decision nodes at the end of the tree. The network would go through the network in order to choose the next word in the sentence based on whether it follows the rules or not. In 2017, the team at DeepMind released a paper detailing the training and performance of their model.
\\[1\baselineskip]
The problems with the model are many. One of the biggest issues is that grammatical rules are ambiguous. For example, the phrase "I came down" could mean I came down from the stairs, or I came down from the sky. In English, there are numerous other examples of sentences that are ambiguous. The model does not have any way to tell the difference between the ambiguous sentences, and so it simply makes a prediction about the entire sentence based on what rules it does know. This makes the model vulnerable to the phenomenon of overfitting. In order to avoid this, models must be trained using a large set of data, so that the model can find useful patterns in the data, and use those patterns to help make predictions. This approach is also prone to errors, as it relies on the training data to help make correct predictions, and the data might be incorrectly labeled, resulting in a model that produces a poor set of predictions. In the case of GPT-2, the model only has about 100 training sentences for each type of phrase (subject-verb-object).
\\[1\baselineskip]
Additionally, languages have different rules for different types of phrasing. In English, the subject-verb-object construction is fairly universal, but other languages have different rules, which makes training a GPT more complicated.
\\[1\baselineskip]
Finally, there is an inherent trade-off between the number of rules and the performance of the model ... 
%As a rule, the more rules the model has, the better it performs. However, as the number of rules increases, the likelihood of the model to overfit the training data increases, and the ability to perform generalizations decreases. For instance, if the model has 1,000 rules ... 
%for an English phrase, it will be able to predict other English phrases with high accuracy. However, it will struggle with generalizing to other languages. If the model has 10,000 rules, it will be able to perform generalizations, but it will be inaccurate with respect to English.
%There have been other grammar learning algorithms that have not yet been released to the public. One such system was used to produce a language translation system that was able to translate into and out of English. However, the system did not learn the rules for the different languages, but instead relied on information about the structure of the language to translate between the two languages. Another system has also been developed by Google DeepMind, which used neural networks to learn grammar. The system was able to learn grammatical rules for two languages, English and French.

}}


\fbox{%
\parbox{0.48\textwidth}{%
\rule{0.48\textwidth}{0.4pt}
{\bf Context}: The 36th International Conference on Machine Learning (ICML 2019) will be held in Long Beach, CA, USA from June 10th to June 15th, 2019. The conference will consist of one day of tutorials (June 10), followed by three days of main conference sessions (June 11-13), followed by two days of workshops (June 14-15).

\rule{0.48\textwidth}{0.4pt}
{\bf Megatron-LM}: With a broad scope, the conference addresses the challenges and opportunities in machine learning for practitioners and researchers. The conference program will feature sessions that cover different machine learning topics, from reinforcement learning to computational biology and from AI to deep learning. It is the world's premier forum for exploring and presenting advances in machine learning.
\\[1\baselineskip]
ICML 2019 will offer an impressive range of sessions, which will be of interest to practitioners, researchers, and students. Some of the topics we will cover include:
\\[1\baselineskip]
Reinforcement learning
\\[1\baselineskip]
Soft Computing
\\[1\baselineskip]
Computational biology
\\[1\baselineskip]
Convolutional networks
\\[1\baselineskip]
Deep learning
\\[1\baselineskip]
Statistical learning theory
\\[1\baselineskip]
Statistical learning machine
\\[1\baselineskip]
Natural language processing
\\[1\baselineskip]
Biology
\\[1\baselineskip]
Awards
\\[1\baselineskip]
Since its first edition, the conference has grown and evolved. Today, it is the premier forum for showcasing state-of-the-art research in machine learning. The conference is widely recognized as the preeminent venue for showcasing the latest research in machine learning, and a place where practitioners, researchers, and students can exchange ideas and insights.
\\[1\baselineskip]
The conference is a peer-reviewed venue that draws its prestige and influence from the excellence of the presentations and the quality of the invited speakers. With over 50\% of all talks selected based on their content by an expert committee, it is clear that the conference's high standards of research and quality keep a very high bar.
%\newline \newline
%For more information about the conference, please visit: ICML.
}}


\fbox{%
\parbox{0.48\textwidth}{%
\rule{0.48\textwidth}{0.4pt}
{\bf Context}: Rescue workers on Tuesday found a flight recorder from the Russian plane that crashed into the Black Sea over the weekend, the defense ministry said. All 92 people aboard the Russian military's Tu-154 plane are believed to have died Sunday morning when it crashed two minutes after taking off from the southern Russian city of Sochi.

\rule{0.48\textwidth}{0.4pt}
{\bf Megatron-LM}: A child holds up a portrait of a dead relative at a makeshift morgue in a building belonging to a local government in the town of Petropavlovsk-Kamchatsky. A Russian Defense Ministry Tu-154 plane with 92 people on board crashed into the Black Sea on Sunday morning shortly after taking off from the southern Russian city of Sochi. (Yuri Kochetkov/EPA)
\\[1\baselineskip]
According to the ministry, the plane was found at a depth of more than 1,000 feet and the recorder was damaged. "The defense ministry has decided to hand over the recorder to the Interstate Aviation Committee, which is the main investigative body in the air crash," the statement said.
\\[1\baselineskip]
The search for the plane's black boxes is being hampered by strong winds and rough seas. Investigators have launched a criminal probe into the crash, which officials said was likely caused by a technical fault.
\\[1\baselineskip]
Russian President Vladimir Putin declared Wednesday a day of mourning for the victims.

}}


\section{Further Scaling Analysis}
\label{app:scaling}

In this section we study the effect of number of attention heads on the scaling results. We also present strong scaling results for our 1.2 billion parameter model.

\subsection{Attention Heads and Scaling} 
This section studies the effect of attention heads on model parallel scaling. To this end, we consider the 8.3 billion parameter configuration with 8-way model parallelism and vary the number of heads from 16 to 32. The results are presented in Table \ref{tab:attn_head_scaling}. As the number of attention heads increases, some of the GEMMS inside the self-attention layer become smaller and also the number of elements in the self attention softmax increases. This results in a slight decrease in scaling efficiency. Future research should be wary of this hyperparameter to design large transformer models that balance model speed and model accuracy.

\begin{table}[hbt!]
\footnotesize
\begin{center}
\caption{Effect of number of attention heads on scaling on 8.3 billion of parameters with 8-way model parallelism.}
\vskip 0.15in
\label{tab:attn_head_scaling}
\begin{tabular}{c|c|c}
\hline \hline
Attention heads & Hidden size per head & Scaling Efficiency \\ \hline
16 & 192 & 82\% \\
24 & 128 & 80\% \\
32 & 96 & 77\% \\ \hline \hline
\end{tabular}
\end{center}
\vskip -0.1in
\end{table}


\subsection{Strong Scaling}

Our model parallelism is primarily designed to enable training models larger than what can fit in the memory of a single GPU, but it can also accelerate the training of smaller models without increasing the batch size. To measure this acceleration we train a model with a fixed 1.2 billion parameters. We use a fixed batch size of 8 samples per iteration and increase the number of GPUs using model parallelism. The results are listed in Table~\ref{tab:strong_scaling}. Using two GPUs makes training $64\%$ faster. Above that we see diminishing returns as the per-GPU computation decreases and the memory bandwidth and communication overheads begin to dominate.

\begin{table}
\centering
\caption{Speedup obtained for the 1.2 billion parameters model using model parallelism while keeping the batch size constant.}
\vskip 0.15in
\begin{tabular}{c|c|c|c|c}
\hline \hline
\# of GPUs & 1 & 2 & 4 & 8 \\
\hline
Speedup & 1.0 & 1.64 & 2.34 & 2.98 \\
\hline \hline
\end{tabular}
\label{tab:strong_scaling}
\vskip -0.1in
\end{table}




\section{Evaluating Language Models Using WikiText103 and LAMBADA}
\label{app:wikilambada}
In this section we detail our evaluation methodology for the WikiText103 dataset \citep{wikitext} and cloze-style prediction accuracy on the LAMBADA dataset\citep{lambada}.

\subsection{Wikitext103 Perplexity}
WikiText103 perplexity is an evaluation criterion that has been well studied over the past few years since the creation of the benchmark dataset. Perplexity is the exponentiation of the average cross entropy of a corpus \citep{perplexity}. This makes it a natural evaluation metric for language models which represent a probability distribution over entire sentences or texts. 

\begin{equation}
 PPL= \exp({-\frac{1}{T_o}\sum_{t}^{T} \text{log} P(t|0:t-1))}
 \label{ppl_eqn}
\end{equation}

To calculate perplexity in (\ref{ppl_eqn}) we tokenize the WikiText103 test corpus according to our subword vocabulary and sum the cross entropy loss from each token $[0, T]$. We then normalize the cross entropy loss by the number of tokens in the original tokenization scheme $T_o$. The WikiText103 test corpus already comes pre-tokenized with word level tokens that prior works have used to compute perplexity. To evaluate our models' perplexities on a level playing field with prior works we must normalize by the original number of tokens, $T_o$, rather than the number of tokens, $T$, actually in the tokenized data fed as input to our model.  This pre-tokenization also introduces artifacts in the text that are not present in our training data. To alleviate this distributional mismatch, we first preprocess the WikiText103 test dataset with invertible detokenizers to remove various artifacts related to punctuation and whitespace. The value of $T_o$ is calculated before this preprocessing. For WikiText103's test set $T_o=245566$ and $T=270329$.

We must also make one further transformer-specific modification to the perplexity calculation. Unlike RNN-based language models, transformers operate on a fixed window input size. Therefore they cannot fully calculate $P(t|0:t-1)$ and can only calculate $P(t|t-w:t-1)$ where $w$ is the size of our context: 1024 tokens. However, calculating this value for every token in our dataset is prohibitively expensive since we must compute approximately $T$ evaluations of a $w$ sized context. To evaluate our models efficiently we take a middle ground approach termed \textit{overlapping evaluation} where we advance the sliding window by some overlap $o$ each time and only compute the cross entropy losses corresponding to the last $o$ tokens of the window. In our experiments we utilize an overlap $o$ of 32, and compute losses over all sliding windows in such a fashion.


\subsection{LAMBADA Cloze Accuracy}
The capability to handle long term contexts is crucial for state of the art language models and is a necessary prerequisite for problems like long-form generation and
document-based question answering. Cloze-style datasets like LAMBADA are designed to measure a model's ability to operate in and reason about these types of long term contexts. Cloze-style reading comprehension uses a context of word tokens $x=x_{1:t}$ with one token $x_j$ masked; the models objective is to correctly predict the value of the missing $j^{\text{th}}$ token. To accurately predict the missing token, the model requires an in-depth understanding of the surrounding context and how language should be used in such a context. LAMBADA uses cloze-style reading comprehension to test generative left-to-right language models by constructing examples of 4-5 sentences where the last word in the context $x_t$ is masked. Our models utilize subword units, so for LAMBADA evaluation we utilize the raw, unprocessed LAMBADA dataset and require that our model predict the multiple subword tokens that make up the word token. We use teacher forcing, and consider an answer correct only when all output predictions are correct. This formulation is equivalent to the original task of word token prediction.



%\bibliography{megatron}
%\bibliographystyle{icml2020}

%\end{document}